\section{Source Code and Benchmark Listings}
\label{app:code}

This appendix provides the complete Python source code used for algorithm execution, automated data generation, and empirical runtime verification. All code is also hosted in the public repository: \url{https://github.com/Harsha-Karimikonda/aoa-project1}.

\subsection{Greedy Scheduler and Empirical Benchmark Harness}
\begin{lstlisting}[language=Python, caption={Greedy Algorithm and Empirical Validation Harness}]
import time, random, gc, csv, platform
from typing import List, NamedTuple
import numpy as np
from scipy.optimize import curve_fit

class Interval(NamedTuple):
    start: float
    finish: float
    request_id: str

def greedy_interval_scheduling(intervals: List[Interval]) -> List[Interval]:
    if not intervals:
        return []
    sorted_intervals = sorted(intervals, key=lambda x: x.finish)
    selected = [sorted_intervals[0]]
    last_finish = sorted_intervals[0].finish
    for current in sorted_intervals[1:]:
        if current.start >= last_finish:
            selected.append(current)
            last_finish = current.finish
    return selected

def generate_random_intervals(n: int, max_time: float = 100000.0, rng=None):
    if rng is None:
        rng = random.Random(42)
    intervals = []
    for i in range(n):
        s = rng.uniform(0, max_time)
        dur = rng.uniform(1.0, 100.0)
        intervals.append(Interval(start=s, finish=s + dur, request_id=f"req_{i}"))
    return intervals

def run_greedy_benchmark():
    n_values = [100, 500, 1000, 5000, 10000, 25000, 50000, 100000, 250000]
    trials = 5
    rng = random.Random(42)
    mean_times = []
    
    for n in n_values:
        trial_times = []
        for _ in range(trials):
            data = generate_random_intervals(n, rng=rng)
            gc.disable()
            t0 = time.perf_counter()
            _ = greedy_interval_scheduling(data)
            t1 = time.perf_counter()
            gc.enable()
            trial_times.append(t1 - t0)
        mean_times.append(float(np.mean(trial_times)))
        
    # Fit theoretical model f(n) = c * n * log2(n)
    def model(n, c):
        return c * n * np.log2(n)
        
    popt, _ = curve_fit(model, np.array(n_values), np.array(mean_times))
    return popt[0], n_values, mean_times
\end{lstlisting}

\subsection{Divide-and-Conquer Closest Pair and Benchmark Harness}
\begin{lstlisting}[language=Python, caption={Divide-and-Conquer Algorithm and Verification Harness}]
import math, time, random, gc
from typing import List, Tuple
import numpy as np

Point = Tuple[float, float]

def euclidean_dist(p1: Point, p2: Point) -> float:
    return math.hypot(p1[0] - p2[0], p1[1] - p2[1])

def brute_force_closest_pair(points: List[Point]):
    n = len(points)
    if n < 2: return float('inf'), None
    min_d, best_pair = float('inf'), None
    for i in range(n):
        for j in range(i + 1, n):
            d = euclidean_dist(points[i], points[j])
            if d < min_d: min_d, best_pair = d, (points[i], points[j])
    return min_d, best_pair

def closest_pair_dnc(points: List[Point]):
    if len(points) < 2: return float('inf'), None
    px = sorted(points, key=lambda p: (p[0], p[1]))
    py = sorted(points, key=lambda p: (p[1], p[0]))
    return _closest_pair_rec(px, py)

def _closest_pair_rec(px: List[Point], py: List[Point]):
    n = len(px)
    if n <= 3: return brute_force_closest_pair(px)
    mid = n // 2
    mid_x = px[mid][0]
    qx, rx = px[:mid], px[mid:]
    left_set = set(qx)
    qy = [p for p in py if p in left_set]
    ry = [p for p in py if p not in left_set]
    
    dist_l, pair_l = _closest_pair_rec(qx, qy)
    dist_r, pair_r = _closest_pair_rec(rx, ry)
    delta, best_pair = (dist_l, pair_l) if dist_l <= dist_r else (dist_r, pair_r)
    
    strip = [p for p in py if abs(p[0] - mid_x) < delta]
    for i in range(len(strip)):
        for j in range(i + 1, min(i + 8, len(strip))):
            if strip[j][1] - strip[i][1] >= delta: break
            d = euclidean_dist(strip[i], strip[j])
            if d < delta: delta, best_pair = d, (strip[i], strip[j])
    return delta, best_pair

def run_dnc_benchmark():
    n_values = [100, 500, 1000, 2500, 5000, 10000, 25000, 50000]
    rng = random.Random(42)
    mean_times = []
    for n in n_values:
        trial_times = []
        for _ in range(5):
            pts = [(rng.uniform(0, 10000), rng.uniform(0, 10000)) for _ in range(n)]
            gc.disable()
            t0 = time.perf_counter()
            _ = closest_pair_dnc(pts)
            t1 = time.perf_counter()
            gc.enable()
            trial_times.append(t1 - t0)
        mean_times.append(float(np.mean(trial_times)))
    return n_values, mean_times
\end{lstlisting}
